\documentclass[a4paper,oneside,article]{memoir}
\usepackage{graphicx} % Required for inserting images
\usepackage[danish]{babel}    % Giver mulighed for dansk orddeling.
\usepackage[utf8x]{inputenc}  % Her sættes tegnsæt til utf8x.
\usepackage{graphicx}	      % Tillader indsættelse af billeder
\usepackage{mathtools}	      % Ekstra matematik.
%\usepackage{url}             % Bruges til at formattere url'er.
\usepackage[colorlinks=true]{hyperref} % Only include hyperref once
%\usepackage[draft]{fixme} % Denne pakke tillader indsættelse af fixme noter ved brug af \fxnote{} som man kan bruge som huske sedler til sig selv.
\usepackage{xcolor}   % Gør så man kan farve tekst
\usepackage{lmodern} % Så man kan compile i TeXworks
\usepackage{alphabeta} % Til at skrive græske bogstaver
\usepackage{subfig} % For at have to billeder ved siden af hinanden
%\usepackage[toc,page]{appendix} % for tilføje bilag
%\usepackage{minted} % Til kodeblok
\usepackage{hyperref} % For clickable TOC (only once)
\usepackage{listings} % Hvis man vil inkludere kode eksempler.
\lstset{language=matlab,breaklines=True} % Syntaks highlight til Matlab.
\usepackage{float} %Mere kontrol over billede-placering
\usepackage{tabularx}


\usepackage{multicol}
\usepackage{microtype} 

\newcommand{\ti}[1]{\cdot 10^{#1}} % Let: gange 10^x
\newcommand{\e}[1]{\emph{e}^{#1}} % Pænere e^x
\newcommand{\vm}[2]{\begin{pmatrix} #1 \\ #2 \end{pmatrix}}
\newcommand{\mat}[1]{\begin{pmatrix} #1 \end{pmatrix}}

\newcommand{\br}[2][1]{%
    \ifcase#1
    \or \left( #2 \right) %1 eller ingenting
    \or \left[ #2 \right] %2
    \or \left\langle #2 \right\rangle %3
    \or \left\{ #2 \right\} %4
    \or \left| #2 \right| %5
    \fi
}

\newcommand{\bil}[1]{
\begin{figure}[H]
    \centering
    \includegraphics[width=0.75\linewidth]{Billeder/#1.png}
    \label{fig:#1}
\end{figure}
}

\newcommand{\pdf}[4]{
\includepdf[pages=#3,scale=.85,pagecommand={\subsection{#2 - \ref{ch:opg}}  \label{asc:#2}}]{pdf/#1.pdf}
\includepdf[pages=\the\numexpr 1 + #3\relax-#4, scale=.85]{pdf/#1.pdf}
}

\newcommand{\h}{\hbar}
\newcommand{\dd}{\partial}
\newcommand{\eq}[1]{\label{eq#1}\tag{#1}}
\renewcommand{\a}{\hat{a}}
\newcommand{\n}{\plainbreak{1}}
\newcommand{\ep}{\epsilon}

\newcommand{\hv}[2][]{%
    \if\relax\detokenize{#1}\relax
        |#2\rangle%
    \else
        \if\relax\detokenize{#2}\relax
            \langle #1|%
        \else
            \langle #1|#2\rangle%
        \fi
    \fi
    }


%____________________________________________________________
\begin{document}


\chapter*{Experimental Physics II, Experiment 1: The Fresnel Relations}
%\chapter*{Foreword}
About fresnel relations and what overall items we used in the lab to investigate this topic, and what results we have. (resume ish)

\begin{multicols}{2}
\chapter{Introduction}
This report covers the Fresnel relations which describe the reflection and transmission of light. For the experiment we use a laser, which we shoot through som lenses and into a semicircular glass, and measure the intensity and angle of the reflected and transmitted light. From this we can use Fresnel's relations to calculate the theoretical intensities: 
\begin{align*}
    R_p &= \frac{\tan^2 (\theta_1 - \theta_2)}{\tan^2 (\theta_1 + \theta_2)}  \eq{7} \\
    T_p &= \frac{\sin 2\theta_1 \sin 2\theta_2}{\sin^2 (\theta_1 + \theta_2) \cos^2 (\theta_1 - \theta_2)}  \eq{8} \\
    R_s &= \frac{\sin^2 (\theta_1 - \theta_2)}{\sin^2 (\theta_1 + \theta_2)}  \eq{9} \\
    T_s &= \frac{\sin 2\theta_1 \sin 2\theta_2}{\sin^2 (\theta_1 + \theta_2)}  \eq{10}
\end{align*}
And we can use Snell's law to calculate the angles 
\begin{align*}
    n_1 \sin\theta_1 = n_2 \sin\theta_2. 
\end{align*}

Where $n_1$ and $n_2$ are the indexes of refraction, which are $n_{\text{air}} = 1.0$ and $n_{\text{glass}} = 1.5$. The angles $θ_1$ and $θ_2$, can be calculated from the measured angles of the disks as seen in \textcolor{red}{figure}.  

\textcolor{gray}{5-6 sentences on the experiment. Describe the purpose and ideas of the project/exercise. Optionally describe the subject of the project/exercise in a broader context. Explain briefly the central physical ideas for the project/exercise. Possibly support the brief explanation with figures. Summarize the experimental methods applied in the investigation.}

\chapter{Experimental setup}
\textcolor{gray}{Describe the applied setup systematically. Explain how measurements are performed. Support the written text with figures/illustrations/photos, for example: a sketch of the experimental setup, photo of the setup taken during the investigation, figures that show direct data (could e.g. be traces obtained with picoscope). The  descriptions  must  be  so  that  a  person  who does  not  know  the  experimental system can repeat the experiments and potentially identify possible flaws. Subsections can often be used to ensure a logical structure of the description.}

\subsection{Equipment}
The following list contains the equipment used, placed in order of appearance in the experimental setup:
\begin{itemize}
    \item Red Diode Laser (PASCO OS-8525A).
    \item Collimating Slits to ensures that the lightbeam remains well-defined and collimated.
    \item Collimator to further refine the beam to maintain parallel rays.
    \item Polarizing Lens to control and adjust the polarization state of the light. (See section \hyperref[sec:cal]{2.2}).
    \item Rotating Disk with Angular Scale, and Brewster's Angle Accessory for measurement of incidence and reflection angles.
    \item Semicircular Dielectric Material (n=1.5).
    \item Focusing Lens.
    \item Transmission Axis Polarization Lens, to focus on either p- or s-polarization. (TAPL)
    \item Optical Filter (10\% Transmission) to prevent detector saturation.
    \item Optical Detector for measuring the intensity of transmitted and reflected light, connected to picoscope for data collection (Saturation at 5V).
\end{itemize}

\subsection{Calibration} \label{sec:cal}
The calibration of the experimental setup is important to ensure accurate measurements of s- and p-polarized light. Since our laser source was mainly emitting p-polarized light, an adjustment was required to attain an equal distribution in intensity of the two polarizations. The calibration itself involves adjustments of a few different components of the setup, and can be a bit tricky to master initially. Since the Detector will saturate at 5V, it was necessary to either let a smaller amount of light through the Collimating slits, by choosing a smaller size or directly by turning down the laser intensity. Then, by making a note of the value of the unpolarized intensity (this can be higher than the saturation value since we cut out all the unwanted polarizations), the appropriate setting of the polarizing lens can be found by alternating between p- and s-polarization with the TAPL, until the intensity is approximately equal for both polarizations (these must be less than the saturation value of 5V). In our case we could find equal intensities at around 4V for each polarization. 

\chapter{Results}
\subsection{Data}
From the experiment we get data for the three measured angles and four intensities. We have data from both glass into air, and air into glass, a part of the data can be seen in \hyperref[tab:data1]{Table 1} for air into glass, and \hyperref[tab:data2]{Table 2} for glass into air. At some points it was impossible to measure the reflected laser, as it was non existent or at almost the same angles as transmitted, this is represented with "-" int the table. 

\subsection{Fitting}
From this we use Snell's law to calculate $θ_2$, which we can use together with the measured $θ_1$ to fit the data to Fresnel's relations. We have plotted the data, the fit and the theoretical values for the normalized intensities of reflected and transmitted light. As fitting parameter we have used the index of refraction for glass, and the fits can be seen in \hyperref[fig:ag]{Figur 1} for air to glass, and in \hyperref[fig:ga]{Figur 2} for glass to air. 

We can see that the measured intensities for reflectance follow the theory very well, and for the fit, we get the fitting parameters shown in table 3


\end{multicols}

\begin{table}[b]
    \centering
    \begin{tabular}{|c|c|c|c|c|c|c|}
        \hline
        $\theta_1$ (°) & $\phi_1$ (°) & $T_{\text{p}} (V)$ & $T_{\text{s}} (V)$ & $\phi_2$ (°) & $R_{\text{p}} (V)$ & $R_{\text{s}} (V)$ \\
        \hline
        0  & 180 & 4.980 & 4.880 & -   & -     & -     \\
        15 & 176 & 5.130 & 4.370 & -   & -     & -     \\
        \hline
        \multicolumn{7}{|c|}{\dots} \\
        \hline
        60 & 164 & 4.541 & 4.539 & 117 & 0.878 & 0.005 \\
        65 & 167 & 3.842 & 4.059 & 128 & 1.136 & 0.0064 \\
        \hline
        \multicolumn{7}{|c|}{\dots} \\
        \hline
        85 & 145 & 0.472 & 0.715 & 168 & 4.638 & 2.339 \\
        90 & 178 & 5.148 & 5.148 & -   & -     & -     \\
        \hline
    \end{tabular}
    \caption{Measured transmission and reflection intensities for different incident angles, for air into glass.}
    \label{tab:data1}
\end{table}

\begin{table}[b]
    \centering
    \begin{tabular}{|c|c|c|c|c|c|c|}
        \hline
        $\theta_1$ (°) & $\phi_1$ (°) & $T_{\text{p}} (V)$ & $T_{\text{s}} (V)$ & $\phi_2$ (°) & $R_{\text{p}} (V)$ & $R_{\text{s}} (V)$ \\
        \hline
        \multicolumn{7}{|c|}{\dots} \\
        \hline
        15 & 172   & 1.848 & 1.785 & 26.5 & 0.114 & 0.085 \\
        20 & 169.5 & 2.181 & 2.332 & 36.5 & 0.117 & 0.073 \\
        \hline
        \multicolumn{7}{|c|}{\dots} \\
        \hline
        42 & 146   & 1.397 & 1.992 & 80.9 & 0.623 & 0.177 \\
        43 & 142.2 & 1.108 & 1.722 & 82.9 & 0.912 & 0.370 \\
        44 & 135.2 & 0.333 & 0.632 & 84.2 & 1.530 & 1.109 \\
        \hline
        \multicolumn{7}{|c|}{\dots} \\
        \hline
        55 & 107   & 1.972 & 2.198 &  -   &  -    &  -    \\
        60 & 112   & 2.252 & 1.841 &  -   &  -    &  -    \\
        \hline
    \end{tabular}
    \caption{Measured transmission and reflection intensities for different incident angles, for glass into air.}
    \label{tab:data2}
\end{table}


\begin{figure}[H]
    \centering
    % First subfigure
    \subfloat[\textcolor{gray}{Reflectance as a function of incident angle. } ]{%
        \includegraphics[width=0.4\textwidth]{Billeder/Rag.png} % Replace with your file
        \label{fig:subfig1}
    }
    \hspace{0.05\textwidth} % Juster afstand mellem billeder
    % Second subfigure
    \subfloat[\textcolor{gray}{Transmission as a function of incident angle. }]{%
        \includegraphics[width=0.4\textwidth]{Billeder/Tag.png}  % Replace with your file
        \label{fig:subfig2}
    }

    % Main figure caption
    \caption{Data, fit and theory for experiment for air into glass. }
    \label{fig:ag}
\end{figure}

\begin{figure}[H]
    \centering
    % First subfigure
    \subfloat[\textcolor{gray}{Reflectance as a function of incident angle. } ]{%
        \includegraphics[width=0.4\textwidth]{Billeder/Rga.png} % Replace with your file
        \label{fig:subfig1}
    }
    \hspace{0.05\textwidth} % Juster afstand mellem billeder
    % Second subfigure
    \subfloat[\textcolor{gray}{Transmission as a function of incident angle. }]{%
        \includegraphics[width=0.4\textwidth]{Billeder/Tga.png}  % Replace with your file
        \label{fig:subfig2}
    }

    % Main figure caption
    \caption{Data, fit and theory for experiment for glass into air. }
    \label{fig:ga}
\end{figure}

\begin{table}[]
\centering
\begin{tabular}{|l|l|l|}
\multicolumn{3}{c}{Air to Glass} \\ \hline
& \text{$n_{glass}$} & \text{Error} \\ \hline
$R_s$ & 1.561 & $\pm 0.060$ \\ \hline
$R_p$ & 1.544 & $\pm 0.556$ \\ \hline
$T_s$ & 1.501 & $\pm 0.080$ \\ \hline
$T_p$ & 1.451 & $\pm 0.735$ \\ \hline
\end{tabular}
\begin{tabular}{|l|l|l|}
\multicolumn{3}{c}{Glass to Air} \\ \hline
& \text{$n_{glass}$} & \text{Error} \\ \hline
$R_s$ & 1.405 & $\pm 0.007$ \\ \hline
$R_p$ & 1.414 & $\pm 0.004$ \\ \hline
$T_s$ & 1.433 & $\pm 0.009$ \\ \hline
$T_p$ & 1.439 & $\pm 0.008$ \\ \hline
\end{tabular}
\caption{Refractive index measurements and errors for Air-to-Glass and Glass-to-Air transitions.}
\label{tab:fit}
\end{table}

\begin{multicols}{2}

\textcolor{gray}{Present measurements and data analyzes from the experimental investigation. Explain how the data has been evaluated, i.e.,what calibrations, calculations, and analyzes (for example fits) have been done. Use figures and tables to clarify the data presentation. Always state uncertainties. For example, the results from a fit is always given with uncertainties and the fitted curve is often plotted together with the data. Apply relevant statistical methods. Reflect critically on the applied methods, i.e., consider for instance the quality of a fit. Consider sources of uncertainty  and inaccuracy.}

\chapter{Discussion}
\textcolor{gray}{Discuss relevant aspects of the investigation. Consider for example the agreement between measured data and theory, reasons for deviations between experiment and theory, possible improvements of the experiment, problems with assumption made in the experiment or theory. Consider using subsections to clarify the discussion. }

\chapter{Conclusion}
\textcolor{gray}{Give a short summary of the idea of the investigation and all relevant conclusions that have been obtained.}

\chapter{References}
\textcolor{gray}{Collect and number all references in order of appearance in the text. Use a consistent format for the references, for example. Authors, Title, Journal name, vol, pages(year). Avoid as far as possible to make direct reference to wikipedia pages and unauthorized webpages, but try instead to find original literature and authorized database.}

\chapter{Appendices}
\textcolor{gray}{Collect other relevant materials that belongs to the obtained data and performed analyzes, for example additional documentation of these.}
\end{multicols}
\end{document}
